\documentclass[a4paper,11pt]{article}
\usepackage[bottom=20mm]{geometry}
\usepackage{mathtools}
% \usepackage{tikz}
% \usetikzlibrary{calc,arrows.meta}

% --- 日本語対応（LuaLaTeX） ---
\usepackage{luatexja}
\usepackage{luatexja-fontspec}
\setmainjfont{HaranoAjiMincho}
\setsansjfont{HaranoAjiGothic}

\title{\vspace{-40mm}Laplace–Beltrami 作用素の導出と\\3次元極座標のラプラシアン}
\author{}

\begin{document}
\maketitle

「3次元極座標のラプラシアン」は皆一度は苦労して通る物理数学の登竜門の１つかと思う。愚直に計算すると結構大変。
そこで、Laplace-Beltrami作用素を使えば、この導出がかなり楽になることはよく知られている。
しかし、そのLaplace–Beltrami作用素自体の導出は意外と骨が折れる。多くの文献では微分幾何などの前提知識が要求される。
ここでは、そういった前提知識なしでLaplace-Beltrami作用素を導出することと、これを使った3次元ラプラシアンの導出を行う。

\section{Laplace–Beltrami 作用素}

一般の座標系 $x^i$について、微小変位の長さの2乗 $ds^2$ は一般には各成分 $dx^i$ の単純な2乗和では書けず、

\begin{equation}
  ds^2 = \mathbf{dx}\cdot\mathbf{dx} = g_{ij}dx^{i}dx^{j}
\end{equation}

\noindent
の形で表される。\footnote{
  (上下)ペアの添字について、$\Sigma$を省略して和を取る記法を\textbf{Einstein規約}という。このpdf内では特に断りのない限りこの記法を用いる。
}
この係数 $g_{ij}$ を\textbf{計量テンソル}という。\\
$g_{ij}$の逆行列を $g^{ij}$、行列式を$g$、$x^i$による偏微分を$\partial_{i}\coloneqq \frac{\partial}{\partial x^{i}}$と書くと、スカラー関数 $\phi$ に作用するラプラシアンは、

\begin{equation}
  \boxed{
    \Delta \phi
    =
    \frac{1}{\sqrt{g}}
    \partial_i
    \left(
    \sqrt{g}\, g^{ij} \partial_j \phi
    \right)
  }
\end{equation}

\noindent
で与えられる。これをLaplace-Beltrami作用素という。

\section{行列式の微分}

後で用いるので、まず行列式の微分公式

\begin{equation}
  \boxed{
    \frac{\partial |A|}{\partial x} = |A|\sum_{i,j} a^{-1}_{ji}\frac{\partial a_{ij}}{\partial x}
    \label{eq:det_diff}
  }
\end{equation}

\noindent
を示しておく。なおここでは一旦Einstein規約を解除しておく。\\
まず、

\begin{equation}
  \frac{\partial{|A|}}{\partial{x}} = \sum_{i,j} \frac{\partial{|A|}}{\partial{a_{ij}}} \frac{\partial{a_{ij}}}{\partial{x}}
  \label{eq:det_diff_01}
\end{equation}

\noindent
となる。ここで、行列$A$の逆行列$A^{-1}$は余因子$\tilde{A}$を用いて、

\begin{equation}
  A^{-1} = \frac{1}{|A|} \tilde{A}
\end{equation}

\noindent
と表わせるのだった。このことから行列式$|A|$は、任意の成分$l$に対して、

\begin{equation}
  |A| = (A\tilde{A})_{ll} = \sum_{k} a_{lk}\tilde{a}_{kl}
\end{equation}

\noindent
と表せられる。

\begin{align}
  \frac{\partial{|A|}}{\partial{a_{ij}}} & = \sum_{k}(\frac{\partial{a_{lk}}}{\partial{a_{ij}}}\tilde{a}_{kl} + a_{lk}\frac{\partial{\tilde{a}_{kl}}}{a_{ij}}) \nonumber            \\
                                         & = \sum_{k} (\frac{\partial{a_{ik}}}{\partial{a_{ij}}}\tilde{a}_{ki} + a_{ik}\frac{\partial{\tilde{a}_{ki}}}{\partial{a_{ij}}}) \nonumber \\
                                         & = |A|a^{-1}_{ji}
\end{align}

\noindent
よって、

\begin{equation}
  \eqref{eq:det_diff_01} = |A|\sum_{i,j} a^{-1}_{ji}\frac{\partial{a_{ij}}}{\partial{x}}
\end{equation}

\noindent
となり、$\eqref{eq:det_diff}$が示せた。

% TODO: CocCommand snippets.openSnippetFiles

\section{Laplace–Beltrami 作用素の導出}

デカルト座標を$\{x^i\}$とすると、ラプラシアンは

\begin{equation}
  \Delta \phi = \delta^{ij} \partial_{i} \partial_{j} \phi
\end{equation}

\noindent
と書ける。一般の座標を$\{x^{i'}\}$として、$\{x^i\}\to\{x^{i'}\}$という座標変換を考える。

\subsection{$\partial_{i}$の座標変換}%


連鎖律から偏微分は、

\begin{equation}
  \boxed{
    \partial_{i} = x^{a'}{}_{,i}\partial_{a'}
  }
\end{equation}

\noindent
と変換される。ここで、

\begin{equation}
  x^{i'}{}_{j} \coloneqq \frac{\partial{x^{i'}}}{\partial{x^{j}}}
\end{equation}

\noindent
である。

\subsection{$g_{i'j'}$の座標変換}%


$ds^2$が不変となる座標変換を考えると、

\begin{equation}
  ds^2 = \delta_{ij}dx^{i}dx^{j} = g_{i'j'}dx^{i'}dx^{j'}
\end{equation}

\noindent
であって、$dx^{i}$の全微分から、

\begin{align}
  \delta_{ij}dx^{i}dx^{j} & = x^{i}{}_{,a'}x^{j}{}_{,b'}\delta_{ij}dx^{a'}dx^{b'} \nonumber \\
                          & = x^{a}{}_{,i'}x^{b}{}_{,j'}\delta_{ab}dx^{i'}dx^{j'}
\end{align}

\noindent

\begin{equation}
  \boxed{
  g_{i'j'} = x^{a}{}_{,i'}x^{b}{}_{,j'}\delta_{ab}
  \label{eq:metric_transform}
  }
\end{equation}

\begin{equation}
  G = {}^t\! X \Delta X
\end{equation}

\begin{align}
  G^{-1} & = X^{-1}\Delta^{-1}\left( {}^t\! X \right)^{-1} \nonumber \\
         & = X^{-1}\Delta^{-1}\,{}^t\! \left( X^{-1} \right)
\end{align}

\begin{equation}
  \boxed{
  g^{i'j'} = x^{i'}{}_{a}x^{j'}{}_{b}\delta^{ab}
  }
\end{equation}

\subsection{ラプラシアンの座標変換}%


\begin{align}
  \Delta \phi & = \delta^{ij} \partial_{i} \partial_{j} \phi \nonumber                                                                                  \\
              & = \partial_{i} \left(\delta^{ij}\partial_{j}\phi\right) \nonumber                                                                       \\
              & = x^{j'}{}_{,i}\partial_{j'}\left(x^{i}{}_{,a'}x^{j}{}_{,b'}g^{a'b'}x^{c'}{}_{,j}\partial_{c'}\phi  \right) \nonumber                   \\
              & = x^{j'}{}_{,i}\partial_{j'}\left(x^{i}{}_{,a'}(x^{c'}{}_{,j}x^{j}{}_{,b'})g^{a'b'}\partial_{c'}\phi  \right) \nonumber                 \\
              & = x^{j'}{}_{,i}\partial_{j'}\left(x^{i}{}_{,a'}\delta^{c'}_{b'}g^{a'b'}\partial_{c'}\phi  \right) \nonumber                             \\
              & = x^{j'}{}_{,i}\partial_{j'}\left(x^{i}{}_{,a'}g^{a'c'}\partial_{c'}\phi  \right) \nonumber                                             \\
              & = x^{j'}{}_{,i}x^{i'}{}_{,a'j'}g^{a'c'}\partial_{c'}\phi + x^{j'}{}_{,i}x^{i}{}_{,a'}\partial_{j'}(g^{a'c'}\partial_{c'}\phi) \nonumber \\
              & = x^{j'}{}_{,i}x^{i'}{}_{,a'j'}g^{a'c'}\partial_{c'}\phi + \delta^{j'}_{a'}\partial_{j'}(g^{a'c'}\partial_{c'}\phi) \nonumber           \\
              & = x^{j'}{}_{,i}x^{i'}{}_{,a'j'}g^{a'c'}\partial_{c'}\phi + \partial_{a'}(g^{a'c'}\partial_{c'}\phi)
  \label{eq:laplace_beltrami_01}
\end{align}

\noindent
ここで、$\eqref{eq:det_diff}$から、

\begin{equation}
  \partial_{i'}g = gg^{j'k'}g_{j'k',i'}
\end{equation}

\noindent
となるので、
\begin{align}
  \frac{\partial_{i'}\sqrt{g}}{\sqrt{g}}
   & = \frac12 g^{a'b'}\,\partial_{i'} g_{a'b'}                     \nonumber \\
   & = \frac12 x^{a'}_{c}\,x^{b'}_{d}\,\delta^{cd}\,
  \partial_{i'}\left( x^{m}_{,a'}\,x^{n}_{,b'}\,\delta_{mn} \right) \nonumber \\
   & = \frac12 x^{a'}_{c}\,x^{b'}_{d}\,\delta^{cd} \left(
  x^{m}_{,a'i'}\,x^{n}_{,b'}\,\delta_{mn}
  + x^{m}_{,a'}\,x^{n}_{,b'i'}\,\delta_{mn}
  \right)                                                           \nonumber \\
   & = \frac12\left(
  x^{a'}_{c}\,x^{m}_{,a'i'}\,x^{n}_{,b'}\,x^{b'}_{d}\,
  \delta^{cd}\,\delta_{mn}
  +
  x^{a'}_{c}\,x^{m}_{,a'}\,x^{n}_{,b'i'}\,x^{b'}_{d}\,
  \delta^{cd}\,\delta_{mn}
  \right)                                                           \nonumber \\
   & = \frac12\left(
  x^{a'}_{c}\,x^{m}_{,a'i'}\,\delta^{c}_{m}
  +
  x^{b'}_{d}\,x^{n}_{,b'i'}\,\delta^{d}_{n}
  \right)                                                           \nonumber \\
   & = \frac12\left(
  x^{a'}_{c}\,x^{c}_{,a'i'}
  +
  x^{b'}_{d}\,x^{d}_{,b'i'}
  \right)                                                           \nonumber \\
   & = x^{a'}_{c}\,x^{c}_{,a'i'}
\end{align}

\begin{align}
  \eqref{eq:laplace_beltrami_01} & = \frac{\partial_{a'}\sqrt{g}}{\sqrt{g}}g^{a'c'}\partial_{c'}\phi + \partial_{a'}(g^{a'c'}\partial_{c'}\phi) \nonumber                            \\
                                 & = \frac{1}{\sqrt{g}} \left( (\partial_{a'}\sqrt{g})g^{a'c'}\partial_{c'}\phi + \sqrt{g}\partial_{a'}(g^{a'c'}\partial_{c'}\phi) \right) \nonumber \\
                                 & = \frac{1}{\sqrt{g}} \partial_{a'} \left(\sqrt{g}g^{a'c'}\partial_{c'}\phi\right)
\end{align}

\noindent
となり、Laplace-Beltrami 作用素を示せた。

\section{3次元極座標のラプラシアンの導出}

\end{document}
